\section*{Chapitre \ref{ch:g12:special-relativity} --- Relativité restreinte~: dilatation du temps}

\begin{solution}{exo:g12:special-relativity:1}
$\gamma = 1/\sqrt{1 - (v/c)^2}$~: $1.01$ ($0.1c$)~; $1.25$ ($0.6c$)~;
$1.67$ ($0.8c$)~; $10.0$ ($0.995c$).
\end{solution}

\begin{solution}{exo:g12:special-relativity:2}
$\gamma = 1.25$, donc le sol mesure $1.25 \times 10.0 =
\qty{12.5}{s}$. Les \qty{10.0}{s} du vaisseau sont le
\omterm{def:g12:special-relativity:proper-time}{temps propre}~: le métronome (une horloge)
est présent aux deux battements.
\end{solution}

\begin{solution}{exo:g12:special-relativity:3}
$v = c\sqrt{1 - 1/\gamma^2}$~: $\gamma = 2$ donne $v = 0.866c$~;
$\gamma = 10$ donne $v = 0.995c$.
\end{solution}

\begin{solution}{exo:g12:special-relativity:4}
$v = \qty{36.1}{m/s}$, $v/c = \num{1.2e-7}$~;
$\gamma - 1 \approx (\num{1.2e-7})^2/2 \approx \num{7.2e-15}$. Une
seconde perdue après $1/\num{7.2e-15} \approx \qty{1.4e14}{s}$ ---
environ quatre millions d'années de conduite.
\end{solution}

\begin{solution}{exo:g12:special-relativity:5}
Le passager roule vers l'éclair avant, donc celui-ci l'atteint le
premier~; par le second postulat les deux éclairs voyagent au même
$c$ sur les demi-distances égales du train \emph{dans son
référentiel}, donc une arrivée plus tôt signifie une frappe plus
tôt. La simultanéité dépend du référentiel~: chaque verdict est
correct dans le sien.
\end{solution}

\begin{solution}{exo:g12:special-relativity:6}
$\gamma = 1/\sqrt{1 - 0.99^2} = 7.09$~; durée de vie labo
$7.09 \times \num{2.6e-8} = \qty{1.8e-7}{s}$.
\omterm{def:g12:projectile-motion:range}{Portée} classique $0.99 \times \num{3.00e8} \times
\num{2.6e-8} \approx \qty{7.7}{m}$~; réelle $7.09 \times 7.7 \approx
\qty{55}{m}$.
\end{solution}

\begin{solution}{exo:g12:special-relativity:7}
$\gamma = 1.5$ à $v = c\sqrt{1 - 1/2.25} \approx 0.75c$~;
$\gamma = 3$ à $\approx 0.94c$. L'asymptote~: $\gamma$ (et le coût
d'\omterm{def:g10:energy-conservation:energy}{énergie} d'une vitesse supplémentaire) diverge
lorsque $v \to c$ --- aucun effort fini n'amène un corps massif à la
vitesse de la lumière.
\end{solution}

\begin{solution}{exo:g12:special-relativity:8}
$\Delta t_0 = 2L/c = 3.0/\num{3.00e8} = \qty{1.0e-8}{s}$~; à $0.8c$,
$\gamma = 5/3$, donc $\Delta t = \qty{1.7e-8}{s}$.
\end{solution}

\begin{solution}{exo:g12:special-relativity:9}
Le \omterm{def:g12:special-relativity:proper-time}{temps propre} appartient à l'horloge
présente aux deux \omterm{def:g12:special-relativity:proper-time}{événements}~: (a)~le
porteur~; (b)~le muon~; (c)~le pilote --- seul le vaisseau est à la
fois au départ et à l'arrivée.
\end{solution}

\begin{solution}{exo:g12:special-relativity:10}
$\gamma - 1 \approx 7700^2/(2 \times \num{9.0e16}) \approx
\num{3.3e-10}$~; sur \qty{1.58e7}{s} l'astronaute retarde de
$\num{3.3e-10} \times \num{1.58e7} \approx \qty{5.2e-3}{s}$ ---
environ cinq millisecondes plus jeune.
\end{solution}

\begin{solution}{exo:g12:special-relativity:11}
$E = \num{1.0e-3} \times (\num{3.00e8})^2 = \qty{9.0e13}{J}$. Une
centrale de \qty{1.0}{GW} livre $\num{1.0e9} \times 86400 \approx
\qty{8.6e13}{J}$ par jour --- un gramme est une journée de sa
production, et $\num{9.0e13}/\num{8.5e5} \approx \num{e8}$~: cent
millions de voitures en vitesse.
\end{solution}

\begin{solution}{exo:g12:special-relativity:12}
Carré et inversion~: $\gamma^2(1 - v^2/c^2) = 1$, donc
$v^2/c^2 = 1 - 1/\gamma^2$, d'où $v = c\sqrt{1 - 1/\gamma^2}$. Pour
$\gamma = 30$~: $v = c\sqrt{1 - 1/900} = 0.99944c$.
\end{solution}

\begin{solution}{exo:g12:special-relativity:13}
Épaisseur contractée $L = \qty{15}{km}/30 = \qty{500}{m}$~; temps de
balayage $500/(0.99944 \times \num{3.00e8}) \approx \qty{1.7e-6}{s} <
\qty{2.2}{\micro s}$. Les deux référentiels s'accordent sur le fait
que le muon atteint le sol --- l'un blâme une durée de vie étirée,
l'autre une atmosphère rétrécie.
\end{solution}

\begin{solution}{exo:g12:special-relativity:14}
Labo~: $3200/\num{3.00e8} \approx \qty{1.1e-5}{s}$ (vitesse
$\approx c$). Référentiel de l'électron~: tube contracté à
$3200/\num{1.0e5} = \qty{3.2}{cm}$, traversé en
$\num{1.1e-5}/\num{1.0e5} \approx \qty{1.1e-10}{s}$.
\end{solution}

\begin{solution}{exo:g12:special-relativity:15}
Terre~: $4.2/0.9 \approx \qty{4.7}{years}$. Avec
$\gamma = 1/\sqrt{1 - 0.81} = 2.29$, l'horloge de la sonde enregistre
$4.7/2.29 \approx \qty{2.0}{years}$.
\end{solution}

\begin{solution}{pb:g12:special-relativity:1}
\textbf{1.} $c\,\Delta t_0 = \num{3.00e8} \times \num{2.2e-6} \approx
\qty{660}{m}$.

\textbf{2.} $\num{15000}/660 \approx 23$~durées de vie.

\textbf{3.} $\mathrm{e}^{-23} \approx \num{e-10}$~: moins d'un muon
sur dix milliards devrait arriver, pourtant le flux est d'un par
\unit{cm^2} par minute --- la physique classique est platement
contredite par le ciel.

\textbf{4.} Durée de vie au sol $30 \times \qty{2.2}{\micro s} =
\qty{66}{\micro s}$~; \omterm{def:g12:projectile-motion:range}{portée} $\approx
\num{3.00e8} \times \num{6.6e-5} \approx \qty{20}{km} > \qty{15}{km}$~:
l'arrivée est expliquée.

\textbf{5.} $v = c\sqrt{1 - 1/900} = 0.99944c$.

\textbf{6.} $\Delta t_0 = 2L/c = 3.0/\num{3.00e8} = \qty{1.0e-8}{s}$.

\textbf{7.} Horizontal $v\,\Delta t/2$~; vertical $L$~; le demi-chemin
oblique mesure $c\,\Delta t/2$ parce que le second postulat fixe la
vitesse du \omterm{def:g12:quantum-world:photon}{photon} à $c$ aussi dans le
\omterm{def:g10:relative-motion:frames}{référentiel sol}.

\textbf{8.} $(c\,\Delta t/2)^2 = L^2 + (v\,\Delta t/2)^2$ donne
$\Delta t^2(c^2 - v^2) = 4L^2$, donc $\Delta t =
(2L/c)/\sqrt{1 - v^2/c^2} = \gamma\,\Delta t_0$.

\textbf{9.} $\gamma = 1/\sqrt{1 - 0.36} = 1.25$~;
$\Delta t = \qty{1.25e-8}{s}$.

\textbf{10.} Celle du vaisseau~: son unique horloge est présente aux
deux tics, donc elle lit le \omterm{def:g12:special-relativity:proper-time}{temps propre}~; le
sol a besoin de deux horloges synchronisées.

\textbf{11.} $\gamma = 1/\sqrt{1 - 0.64} = 1/0.6 = 5/3 \approx 1.67$.

\textbf{12.} $10.0/(5/3) = \qty{6.0}{years}$.

\textbf{13.} Aller $\qty{5.0}{years}$ à $0.8c$~: $4.0$
\omterm{def:g10:orders-of-magnitude:lightyear}{années-lumière}.

\textbf{14.} $10.0 - 6.0 = \qty{4.0}{years}$~: l'astronaute est
quatre ans plus jeune que son jumeau.

\textbf{15.} La situation n'est pas symétrique~: seule l'astronaute
fait demi-tour --- elle accélère et change de
\omterm{def:g12:special-relativity:inertial}{référentiel galiléen}, tandis que le jumeau
terrestre reste dans un seul --- donc seule son horloge enregistre le
temps le plus court.

\textbf{16.} $v = 2\pi \times \num{2.66e7}/\num{43200} \approx
\qty{3.9e3}{m/s}$.

\textbf{17.} $\gamma - 1 \approx (\num{3.87e3})^2/(2 \times
\num{9.0e16}) \approx \num{8.3e-11}$.

\textbf{18.} $\num{8.3e-11} \times \num{86400} \approx \qty{7.2e-6}{s}
= \qty{7.2}{\micro s}$ par jour de retard.

\textbf{19.} $c \times \num{7.2e-6} \approx \qty{2.2}{km}$ d'erreur de
position par jour.

\textbf{20.} Dérive nette $45 - 7.2 \approx \qty{38}{\micro s}$ par
jour d'avance~; erreur $c \times \num{38e-6} \approx \qty{11}{km}$
par jour --- la facture relativiste que chaque horloge GPS paie
d'avance, désaccordée au sol pour tic-tac juste en
\omterm{def:g10:universal-gravitation:satellite}{orbite}.
\end{solution}
